\documentclass[10pt,a4paper]{amsart}
\usepackage[margin=19mm]{geometry}
\usepackage{amsmath,amssymb,mathtools}
\usepackage[scr=rsfs,cal=euler]{mathalfa}
\usepackage{xfrac}
\usepackage[unicode,hidelinks,bookmarks=false]{hyperref}
\newcommand{\ie}{i.e.\ }
\newcommand{\cf}{cf.\ }
\newcommand{\resp}{respectively\ }
\newcommand{\Subsection}{\subsection}
\providecommand{\nobreakdash}{\nobreak\hbox{-}}
\setlength{\emergencystretch}{2em}
\multlinegap0pt
\usepackage{bm}
\usepackage{bbm}
\usepackage{dsfont}
\usepackage{xspace}
\usepackage{cancel}
\usepackage{mathtools}
\usepackage{amsthm}
\makeatletter
\@ifundefined{theorem}{\newtheorem{theorem}{Theorem}}{}
\@ifundefined{lemma}{\newtheorem{lemma}[theorem]{Lemma}}{}
\@ifundefined{proposition}{\newtheorem{proposition}[theorem]{Proposition}}{}
\makeatother
\newcommand{\TT}{\mathcal{T}}
\title[$\sigma$-Irregularity of Trees with Prescribed Maximum Degre]{$\sigma$-Irregularity of Trees with Prescribed Maximum Degree: A Majorization--Duality--Stability Framework}
\author{Jasem Hamoud, Duaa Abdullah}
\date{Source: arXiv:2608.12991v1 (CC BY 4.0)}
\begin{document}
\maketitle

\noindent\emph{Source and licence.} $\sigma$-Irregularity of Trees with Prescribed Maximum Degree: A Majorization--Duality--Stability Framework. Version arXiv:2608.12991v1 (CC BY 4.0), distributed under the Creative Commons Attribution 4.0 International licence, as recorded in the arXiv licence metadata for this exact version and shown on the abstract page https://arxiv.org/abs/2608.12991v1. Full rights notice: licenses/arxiv-2608.12991-CC-BY-4.0.txt.\par\smallskip

\noindent\textit{Editorial scope.} Counted unit: the Theorem whose source label is \texttt{thm001internal} in the article (source0.tex lines 645--647, source label \texttt{thm001internal}), delivered together with the article's own complete original proof of it (source0.tex lines 648--771, one \texttt{proof} environment of 124 lines). No mathematics is written, repaired or re-derived: statement and proof are the article's own text line for line. Required context retained uncounted: prop001pathmonotonicity (lines 482--488); lem001atmostoneinternalleaf (lines 577--579). Every definition, display and label of those lines is retained unchanged. Omitted from this package and not redistributed: 1--481, 489--576, 580--644, 772--1503. Nothing inside a retained excerpt is omitted or altered: every retained body is delivered exactly as the article's lines are written, so no omission receipt of any kind accompanies this package. Citations: the retained excerpts carry no citation command at all, so no bibliography entry of the article is transcribed and no reference is redistributed. Reference adaptation: none was needed and none was made; every reference inside the delivered unit resolves inside it, so no unresolved reference or citation is produced. Printed slips kept literal: no printed word, symbol, hypothesis or number of the counted statement or of its proof was altered. Layout and class: the article's own class is \texttt{amsart} with packages including fontenc, geometry, graphicx, amssymb, stmaryrd, enumitem, tikz, float, url, mathtools, amsmath, amsfonts, amsthm, mathrsfs; the wrapper is the lane's pinned \texttt{amsart} editorial wrapper at 10pt on A4, which restates the theorem-like environments the retained text needs and adds the article's own macro definitions that the retained text needs (\texttt{\textbackslash TT}), with extra wrapper packages (bm, bbm, dsfont, xspace, cancel, mathtools). The delivered numbering is the unit's own. Assets: no retained span contains a figure environment, a graphics-inclusion command, a PDF-page inclusion, a table or a TikZ picture; no image file is copied into this package, the package declares no asset dependency and no image file is redistributed with it. Licence: version arXiv:2608.12991v1 of this article is distributed under the Creative Commons Attribution 4.0 International licence, as recorded in the arXiv licence metadata for this exact version and shown on the abstract page https://arxiv.org/abs/2608.12991v1. A full rights notice is delivered at licenses/arxiv-2608.12991-CC-BY-4.0.txt. Characters: every retained excerpt is pure ASCII, so the delivered engine, XeLaTeX with its default Latin Modern text font, reports no missing character.
\begin{proposition}~\label{prop001pathmonotonicity}
Let $T\in \TT_{n(i,j),\,\Delta}$ be a maximal tree,
$\mathcal P=u\,x\cdots y\,v$ be a path in $T$ where $x=y$ is allowed if
$d_T(u,v)=2$, and $x,y$ may coincide with $u,v$ if $d_T(u,v)\leqslant 2$. If
$d(u)>d(v)$, then $d(x)\leqslant d(y)$. Symmetrically, if $d(x)>d(y)$, then
$d(u)\leqslant d(v)$.
\end{proposition}

\begin{lemma}~\label{lem001atmostoneinternalleaf}
If $n\geqslant \Delta+4$, then at most one internal leaf of $T\in \TT_{n(i,j),\,\Delta}$ has degree strictly less than $\Delta$.
\end{lemma}

\begin{theorem}~\label{thm001internal}
Let $T\in \TT_{n(i,j),\,\Delta}$ be a maximal tree with maximum degree $\Delta$. If $n\geqslant 7$ for $\Delta=3$, or $n\geqslant 2\,\Delta$ for $\Delta\geqslant 4$. Then, every internal leaf of $T$ has degree exactly $\Delta$.
\end{theorem}

\begin{proof}
Suppose toward a contradiction that some
internal leaf $u$ of $T$ has $d(u)<\Delta$. By
Lemma~\ref{lem001atmostoneinternalleaf}, every other internal leaf of $T$
has degree exactly $\Delta$. Since $n\geqslant 2\,\Delta$ and $d(u)<\Delta$, $T$
contains at least one core vertex. Indeed, if $T$ had no core vertex, its non leaf vertices would be exactly its internal leaves, contributing at most
$n-(\text{leaves}) \leqslant \Delta$ vertices with total degree at most $\Delta^2$
edges attaching leaves plus $O(\Delta)$ internal edges, which is insufficient
to reach $n\geqslant 2\Delta$ together with $d(u)<\Delta$.

Let $v$ be any core vertex, and let $w$ be a neighbor of $v$ with $d_T(w,u)>d_T(v,u)$ and $w$ not a leaf such $w$ exists since $v$, being a core vertex, is not itself an internal leaf. Thus, it has at least two non leaf neighbors, at least one of which is farther from
$u$ than $v$ is). If $d(v)>d(u)$, according to Proposition~\ref{prop001pathmonotonicity}
applied to the path from $u$ through $v$ to $w$ where $d(w)\geqslant 2$,
contradicting that $w$ is not a leaf. Hence $d(v)\leqslant d(u)$ for every core
vertex $v$.

We now distinguish two cases according to the maximum number of non-leaf
neighbors of any core vertex.

\medskip

\noindent \textbf{Case 1.} \textit{Every core vertex has exactly two non-leaf neighbors.}

\smallskip

\textbf{Subcase~1.1.} \emph{Every core vertex has degree $2$.} Then $T$ has exactly
two internal leaves, $u$ and some $v$ with $d(v)=\Delta$ by considering 
Lemma~\ref{lem001atmostoneinternalleaf}, joined by a path $\mathcal P$ whose
interior vertices are precisely the core vertices, all of degree $2$.

If $\mathcal P$ has exactly one interior core vertex $w$. Then, $n=1+(d(u)-1)+1+1+(\Delta-1)+\cdots$.  Thus, by counting all vertices of $T$ gives
$n=d(u)+1+\Delta$. Since $n\geqslant 2\,\Delta$ and $d(w)=2$ one
obtains $d(u)=\Delta-1$.  Let $T'\in \TT_{n(i,j),\,\Delta}$ where  $T' = T-vw+vu$. Then, the edges $vw,uv$ contribute
$0-(\Delta-2)^2$ to $\sigma(T')-\sigma(T)$ where edge $vw$, contribution $0$ since
$d(v)=d(w)=\ldots$, the edge $uw$   contributes $(\Delta-1)^2-(\Delta-3)^2$, and the $\Delta-2$ leaves attached to $u$ contribute $(\Delta-1)^2-(\Delta-2)^2$ each. Therefore,
\begin{equation}~\label{eqq01subcase1a}
\sigma(T')-\sigma(T)=\Delta^2+\Delta-6,
\end{equation}
which is strictly positive for every $\Delta\geqslant 3$. Indeed
$\Delta^2+\Delta-6=(\Delta+3)(\Delta-2)>0$ for $\Delta>2$, a contradiction.

\medskip 

If $\mathcal P$ has more than one interior vertex, let $\mathcal P'=zwu$ be the subpath nearest $u$ where $w,z$ are core vertices of degree $2$, or $z=v$ if $\mathcal P$ has exactly two interior vertices. Let $T'\in \TT_{n(i,j),\,\Delta}$ where $T'=T-zw+zu$. The edges $zw,zu$ contribute $(d(u)-2)^2$, the edge $wu$ contributes $(d(u)-1)^2-(d(u)-2)^2$, and the $d(u)-1$ leaves attached to $u$ contribute $d(u)^2-(d(u)-1)^2$ each. Thus, 
\begin{equation}~\label{eqq01subcase1amulticore}
\sigma(T')-\sigma(T)\! =\! (d(u)\!-\!2)^2\!+\!(d(u)\!-\!1)^2\!-\!(d(u)\!-\!2)^2\!+\!(d(u)\!-\!1)\big[d(u)^2\!-\!(d(u)\!-\!1)^2\big], 
\end{equation}
which is holds $\sigma(T')-\sigma(T)=3d(u)^2-5d(u)+2,$ strictly positive for every $d(u)\geqslant 2$, it factors as $(3d(u)-2)(d(u)-1)$, both factors positive for $d(u)\geqslant 2$, a contradiction.

\textbf{Subcase~1.2.} \emph{ Some core vertex $v$ has degree $d(v)\geqslant 3$.} Assume that  $v$ to be the degree $d(v)\geqslant 3$ core vertex closest to $u$. Thus, every interior vertex of  $\mathcal P$ from $u$ to $v$ has degree $2$. Since $v$ has exactly two non leaf
neighbors   and $d(v)\geqslant 3$, $v$ has at least one pendant leaf $v_1$. Let $T'\in \TT_{n(i,j),\,\Delta}$ where $T'=T-v_1v+v_1u$. Since $v$ is a core vertex, $d(v)\leqslant d(u)$, the pair $v_1v,v_1u$ contributes strictly positively to $\sigma(T')-\sigma(T)$, and the edges of $\mathcal P$ likewise contribute a strictly positive sum where each edge of $\mathcal P$ sees one endpoint's ``target'' degree effectively shift from $d(v)$ toward
$d(u)$, and $d(v)\leqslant d(u)$ makes each such shift favorable.

Let $v_2$ be the non leaf neighbor of $v$ not on $\mathcal P$, and $u_1$ a leaf
attached to $u$. If $d(v_2)\geqslant d(v)$, the edge $v_2v$ contributes positively;
otherwise its contribution is bounded below by $(d(v)-2)^2-(d(v)-1)^2$. Hence, $\sigma_1$ of $v_2v$ and $u_1u$ satisfies
\begin{equation}~\label{eqq01subcase1bsigma1}
\sigma_1 \geqslant (d(v)-2)^2-(d(v)-1)^2 + d(u)^2-(d(u)-1)^2. 
\end{equation}
Eq.~\eqref{eqq01subcase1bsigma1} emphasizes that $\sigma_1 \geqslant 2d(u)-2d(v)+2 > 0$. Let $\mathcal L_v$ (resp.\ $\mathcal L_u$) be the leaf edges at $v$
(resp.\ $u$) other than $v_1v$ (resp.\ $u_1u$). Since $d(v)\leqslant d(u)$ and $v$
has one fewer available leaf slot than $u$ when $v_2v$ occupies a non-leaf
slot at $v$), $|\mathcal L_v|<|\mathcal L_u|$. According to~\eqref{eqq01subcase1amulticore},~\eqref{eqq01subcase1bsigma1} and
the positive contribution of $\mathcal P$'s edges,
\begin{equation}~\label{eqq01subcase1bfinal}
\sigma(T')-\sigma(T) > |\mathcal L_v|\Big[(d(v)-2)^2-(d(v)-1)^2\Big]
- |\mathcal L_u|\Big[d(u)^2-(d(u)-1)^2\Big]
> |\mathcal L_v|\big(2d(u)-2d(v)+2\big),
\end{equation}
where $\sigma(T')-\sigma(T) >0$ a contradiction.

\medskip

\noindent \textbf{Case 2.} \textit{Some core vertex $v$ has at least three non leaf neighbors.} Assume that $v$ to be such a vertex closest to $u$, and let $\mathcal P$ be the path from $u$ to $v$, all of whose interior vertices have degree $2$.

\smallskip

\textbf{Subcase~2.1.} \emph{Every non leaf neighbor of $v$ other than possibly the one
on $\mathcal P$ is an internal leaf.}  Then $v$ is the unique core vertex of $T$ and $u$
is adjacent to $v$. If $v$ has a pendant leaf $w$, let $T'\in \TT_{n(i,j),\,\Delta}$ where $T'=T-wv+wu$; since
$d(v)\leqslant d(u)$, the pair $wv,wu$ and the edge $uv$ both contribute positively,
each internal leaf neighbor $z$ of $v$   contributes positively via $vz$,
and only the leaf edges $\mathcal L_v$ at $v$  can contribute
negatively, each by at least $(d(v)-2)^2-(d(v)-1)^2$; since
$|\mathcal L_v|<|\mathcal L_u|$, the same telescoping as
\eqref{eqq01subcase1bfinal} gives $\sigma(T')-\sigma(T)>0$, a contradiction.
If instead $v$ has no pendant leaf, let $z$ be any leaf of $T$ not attached to
$u$ and consider $T'=T-uv+uz$. Then, 
\begin{equation}~\label{eqq01subcase2anoleaf}
\sigma(T')-\sigma(T) >
(d(v)-1)\Big[(\Delta\!-\!(d(v)\!-\!1))^2\!-\!(\Delta\!-\!d(v))^2\Big]
\!+\!(\Delta\!-\!2)^2\!-\!(\Delta\!-\!1)^2.
\end{equation}
Thus, $\sigma(T')-\sigma(T) >(d(v)-2)(2\Delta-2d(v)-1) > 0,$
since $3\leqslant d(v)<\Delta$ when $v$ is a core vertex with $d(v)\leqslant d(u)<\Delta$,
 by using $d(u)<\Delta$, a contradiction.

\smallskip

\textbf{Subcase~2.2.} \emph{If $v$ has a non leaf neighbor that is itself a core vertex.} 
Let $u_1,v_1$ be the neighbors of $u,v$ on $P$. If $v_1$ is the only
core vertex neighbor of $v$, then $v$'s remaining non-leaf neighbor $v_2$ is
an internal leaf, so $d(v_2)=\Delta$ according to
Lemma~\ref{lem001atmostoneinternalleaf}. Since $d(u)<\Delta=d(v_2)$,
Proposition~\ref{prop001pathmonotonicity} applied to the path from $u$ through
$u_1$ to $v_2$   forces $d(u_1)\geqslant d(v)$, contradicting
$d(u_1)\leqslant d(v)$.

Hence $v$ has a core vertex neighbor $v_2\ne v_1$. Then, assume that of minimum degree
among all such neighbors where $d(v_2)\leqslant d(u)$.
Let $v_3\ne v_1$ be the third non leaf neighbor of $v$ where minimality of $v_2$
gives $d(v_2)\le d(v_3)$. Let $T'=T-v_2v+v_2u$ and the edges of $\mathcal P$ contribute a
strictly positive sum  as $d(v)\leqslant d(u)$). In this case, let $u_2,u_3$ be two leaves at $u$, and let $\sigma_1$ be the combined contribution of $v_2v, v_2u, v_3v, u_2u,$ and $u_3u$ gives
\begin{equation}~\label{eqq001subcase2bsigma1}
\sigma_1 = \big(d(u)-d(v_2)\big)^2+3d(u)+d(v)\big(d(v)-2\big)
+3\big(d(u)-d(v_2)\big)+2d(v_3),
\end{equation}
where $\sigma_1>0$ and by using $d(u)\geqslant d(v_2)$. With $\mathcal L_v,\mathcal L_u$ the remaining leaf edges at $v,u$
respectively ($|\mathcal L_v|=d(v)-3\leqslant d(u)-3=|\mathcal L_u|$ since $d(v)\leqslant d(u)$), which \eqref{eqq01subcase1bfinal} yields
$\sigma(T')-\sigma(T)>0$, a contradiction.

In every case we reach a contradiction, so no internal leaf of $T$ can have
degree less than $\Delta$.
\end{proof}

\end{document}
